\documentclass[UTF8]{ctexart}
\usepackage{xeCJK}
\usepackage{geometry}
\geometry{left=2cm,right=2cm,top=2.5cm,bottom=2.5cm}
\setlength{\headheight}{12.7pt}

\usepackage{titlesec}
\titleformat{\section}
  {\normalfont\large\bfseries}
  {\thesection}
  {1em}
  {}
\titlespacing*{\section}
  {0pt}
  {3.5ex plus 1ex minus .2ex}
  {2.3ex plus .2ex}

\usepackage{amsmath}
\usepackage{amssymb}
\usepackage{amsthm}
\usepackage{enumitem}
\setlist[enumerate]{itemsep=0pt,topsep=0pt,parsep=0pt,leftmargin=0pt,itemindent=2.5em}
\setlist[itemize]{itemsep=0pt,topsep=0pt,parsep=0pt,leftmargin=0pt,itemindent=2.5em}
\usepackage{fancyhdr}
\pagestyle{fancy}
\fancyhf{}
\usepackage{graphicx}
\usepackage{hyperref}
\usepackage{indentfirst}
\usepackage{lmodern}

\begin{document}
\setlength{\abovedisplayskip}{1pt}
\setlength{\belowdisplayskip}{1pt}

\section{一阶紧致性定理}

\hypertarget{sec:定理1.1}{}
\noindent\textbf{定理1.1 (一阶紧致性定理)}

给定一阶语言$\mathcal{L}$. 对任意一族公式$\Gamma$,
如果$\Gamma$的任意有限子集都一致, 那么$\Gamma$一致.

\noindent{\kaishu 证明.}

假设$\Gamma$的任意有限子集都一致, 定义
\[
    I=\{i\subset\Gamma\mid i\,\text{有限}\},
\]
则对任意的$i\in I$, 存在$\mathcal{L}$-环境$(A_i,\sigma_i)$满足$i$.

对任意的$\varphi\in\Gamma$, 定义$\Omega_\varphi=\{i\in I\mid i\ni\varphi\}$
(即$\Gamma$所有含$\varphi$的有限子集), 考虑全体$\Omega_\varphi$组成的集合
\[
    B=\{\Omega_\varphi\mid\varphi\in\Gamma\}\subset\mathcal{P}(I).
\]
任取它的非空有限子集$\{\Omega_{\varphi_1},\cdots,\Omega_{\varphi_n}\}\subset B$,
注意到
\[
    \{\varphi_1,\cdots,\varphi_n\}\in
    \Omega_{\varphi_1}\cap\cdots\cap\Omega_{\varphi_n},
\]
即$B$的非空有限子集的交集都非空, 从而可以扩张成$I$上的一个超滤$\mathcal{U}\supset B$.

\vspace{3pt}
最后, 考虑超积结构$^{\displaystyle P}\!/_{\displaystyle\mathcal{U}}$和它的变元赋值
\[
    ^{\displaystyle S}\!/_{\displaystyle\mathcal{U}}:
    \mathcal{V}\to\,^{\displaystyle P}\!/_{\displaystyle\mathcal{U}},\,
    x\mapsto[(\sigma_i(x))_{i\in I}]\\[3pt]
\]
组成的环境$\left(\,^{\displaystyle P}\!/_{\displaystyle\mathcal{U}},\,
^{\displaystyle S}\!/_{\displaystyle\mathcal{U}}\right)$. 因为对任意的
$\varphi\in\Gamma$, 存在$\Omega_\varphi\in\mathcal{U}$使得对其中任意的$i$有
\[
    i\in\Omega_\varphi\implies i\ni\varphi\implies
    (A_i,\sigma_i)\models\varphi,\\[3pt]
\]
于是根据 \href{02 Łoś 定理#nameddest=sec:定理1.1}{Łoś定理}可知
$\left(\,^{\displaystyle P}\!/_{\displaystyle\mathcal{U}},\,^{\displaystyle S}
\!/_{\displaystyle\mathcal{U}}\right)\models\varphi$,
从而该环境满足$\Gamma$. \hfill $\qedsymbol$

\vspace{20pt}

\hypertarget{sec:定理1.2}{}
\noindent\textbf{定理1.2 (一阶紧致性定理)}

给定一阶语言$\mathcal{L}$. 对任意一族公式$\Gamma$和公式$\varphi$,
如果$\Gamma\models\varphi$, 那么存在$\Gamma$的有限子集$\Delta$
使得$\Delta\models\varphi$.

\noindent{\kaishu 证明.}

$\Gamma\cup\{\neg\varphi\}$不一致, 由上述定理可知, 存在有限的
$S\subset\Gamma\cup\{\neg\varphi\}$不一致.

令$D=S\cup\{\neg\varphi\}$, 则$D$也是$\Gamma\cup\{\neg\varphi\}$的有限子集
并且也不一致(因为$S$不一致).

再令$\Delta=D\setminus\{\neg\varphi\}$, 则$\Delta$是$\Gamma$的有限子集.
此时$\Delta\cup\{\neg\varphi\}=D$不一致, 所以$\Delta\models\varphi$.
\hfill $\qedsymbol$

\end{document}
